\documentclass[aspectratio=169, 11pt]{beamer}

% ------------------------------------------------
% PACOTES
% ------------------------------------------------

\usepackage[utf8]{inputenc}
\usepackage[T1]{fontenc}
\usepackage[brazil]{babel}
\usepackage{listings}
\usepackage{xcolor}
\usepackage{booktabs}
\usepackage{graphicx}

% ------------------------------------------------
% IDENTIDADE VISUAL
% ------------------------------------------------

\definecolor{background}{RGB}{20,25,35}
\definecolor{primary}{RGB}{90,170,255}
\definecolor{secondary}{RGB}{130,210,180}
\definecolor{codebg}{RGB}{30,36,48}
\definecolor{textcolor}{RGB}{230,235,245}

\usetheme{Madrid}

\setbeamercolor{normal text}{
    fg=textcolor,
    bg=background
}

\setbeamercolor{frametitle}{
    fg=primary,
    bg=background
}

\setbeamercolor{title}{
    fg=primary
}

\setbeamercolor{structure}{
    fg=secondary
}

\setbeamercolor{block title}{
    fg=background,
    bg=primary
}

\setbeamercolor{block body}{
    fg=textcolor,
    bg=codebg
}

\setbeamertemplate{navigation symbols}{}

% ------------------------------------------------
% CONFIGURAÇÃO DO CÓDIGO
% ------------------------------------------------

\lstdefinestyle{pythonstyle}{
    language=Python,
    backgroundcolor=\color{codebg},
    basicstyle=\ttfamily\small\color{textcolor},
    keywordstyle=\color{primary}\bfseries,
    stringstyle=\color{secondary},
    commentstyle=\color{gray}\itshape,
    numbers=left,
    numberstyle=\tiny\color{gray},
    breaklines=true,
    frame=single,
    rulecolor=\color{primary},
    showstringspaces=false,
    tabsize=4
}

% ------------------------------------------------
% INFORMAÇÕES DO DOCUMENTO
% ------------------------------------------------

\title{Analisador Léxico}
\subtitle{Construção de um Scanner para a linguagem MiniLang}

\author{Samuel}
\institute{Ciência da Computação}

\date{\today}

% ------------------------------------------------
% DOCUMENTO
% ------------------------------------------------

\begin{document}

% CAPA
\begin{frame}
    \titlepage
\end{frame}

% SUMÁRIO
\begin{frame}{Sumário}
    \tableofcontents
\end{frame}

% ------------------------------------------------
\section{Introdução}

\begin{frame}{Introdução}

    A análise léxica representa a primeira etapa do
    processo de compilação de uma linguagem de programação.

    \vspace{0.4cm}

    Seu objetivo consiste em percorrer o código-fonte,
    identificando e classificando sequências de caracteres
    em unidades denominadas \textbf{tokens}.

    \vspace{0.4cm}

    \begin{block}{Proposta do projeto}
        Desenvolver um analisador léxico utilizando Python
        e expressões regulares para reconhecer os elementos
        fundamentais da linguagem MiniLang.
    \end{block}

\end{frame}

% ------------------------------------------------
\section{Definição Formal}

\begin{frame}{Definição Formal}

	A utilização do Autômato Finito Determinístico (AFD) 
	justifica-se pela sua capacidade de reconhecer padrões 
	de caracteres de maneira organizada e previsível. 

    Um Autômato Finito Determinístico é representado pela tupla:
    
    \begin{equation}
		M = (Q, \Sigma, \delta, q_0, F)
	\end{equation}
    
    \begin{table}[h]
    \centering
    \caption{Elementos da tupla matemática do AFD}
    \begin{tabular}{|c|l|}
        \hline
        \textbf{Elemento} & \textbf{Significado} \\ \hline
        $Q$ & Conjunto de estados do autômato. \\ \hline
        $\Sigma$ & Alfabeto de entrada (caracteres aceitos). \\ \hline
        $\delta$ & Função de transição entre os estados. \\ \hline
        $q_0$ & Estado inicial. \\ \hline
        $F$ & Conjunto de estados finais (aceitação). \\ \hline
    \end{tabular}
    \label{tab:elementos-afd}
	\end{table}

\end{frame}


\begin{frame}
	
	
	No contexto deste projeto, o autômato será utilizado como base para o reconhecimento dos padrões 			presentes na MiniLang. A partir da leitura dos caracteres de entrada, o analisador identifica as 			sequências correspondentes a cada categoria léxica, classificando-as em tokens, como palavras-				chave, identificadores, literais numéricos e operadores aritméticos.
	
   \begin{itemize}
        \item \textbf{Palavras-chave:} if, else, while, int, float.
        \item \textbf{Identificadores:} contador, idade, valor.
        \item \textbf{Números inteiros:} 10, 25, 100.
        \item \textbf{Números decimais:} 10.5, 3.14.
        \item \textbf{Operadores aritméticos:} +, -, *, /.
    \end{itemize}

\end{frame}

% ------------------------------------------------
\section{Diagrama de Transições}

\begin{frame}{Diagrama de Transições}

    Tokens são unidades léxicas reconhecidas durante
    a leitura do código-fonte.

    \vspace{0.3cm}

    Considere a seguinte instrução:

    \begin{block}{Exemplo}
        \texttt{int idade = 20;}
    \end{block}

    O analisador identifica:

    \begin{center}
    \begin{tabular}{ll}
        \toprule
        \textbf{Token} & \textbf{Classificação} \\
        \midrule
        int & Palavra-chave \\
        idade & Identificador \\
        = & Operador \\
        20 & Número inteiro \\
        ; & Delimitador \\
        \bottomrule
    \end{tabular}
    \end{center}

\end{frame}

% ------------------------------------------------
\section{Minimização (se aplicável)}

\begin{frame}[fragile]{Minimização (se aplicável)}

    \begin{itemize}
        \item \textbf{Python:} linguagem utilizada na implementação;
        \item \textbf{re:} biblioteca para expressões regulares;
        \item \textbf{Regex:} mecanismo de reconhecimento dos padrões;
        \item \textbf{MiniLang:} linguagem utilizada como referência.
    \end{itemize}

\end{frame}

\begin{frame}[fragile]{Definição dos tokens}

\begin{lstlisting}[style=pythonstyle]
PALAVRAS_CHAVE = {
    "int", "float", "if",
    "else", "while", "print"
}

REGRAS = [
    ("NUM_FLOAT", r"\d+\.\d+"),
    ("NUMINT", r"\d+"),
    ("IDENTIFICADOR",
     r"[a-zA-Z][a-zA-Z0-9_]*"),
    ("OPERADOR", r"==|!=|>=|<=|[+\-*/=<>]"),
    ("DELIMITADOR", r"[(){},;]"),
    ("ESPACO", r"\s+"),
]
\end{lstlisting}

\end{frame}

\begin{frame}[fragile]{Construção do padrão}

\begin{lstlisting}[style=pythonstyle]
import re

PADRAO = re.compile(
    "|".join(
        f"(?P<{nome}>{regra})"
        for nome, regra in REGRAS
    )
)
\end{lstlisting}

    \vspace{0.3cm}

    Cada expressão regular é associada a um grupo nomeado,
    permitindo identificar qual regra reconheceu determinado
    trecho do código-fonte.

\end{frame}

% ------------------------------------------------
\section{Implementação de Apoio}

\begin{frame}{Implementação de Apoio}

    \begin{enumerate}
        \item Receber o código-fonte;
        \item Percorrer os caracteres sequencialmente;
        \item Aplicar as expressões regulares;
        \item Identificar o tipo de token;
        \item Descartar espaços em branco;
        \item Verificar palavras reservadas;
        \item Armazenar os tokens identificados;
        \item Reportar caracteres inválidos.
    \end{enumerate}

\end{frame}

\begin{frame}[fragile]{Exemplo de entrada}

\begin{lstlisting}[style=pythonstyle]
codigo = """
int idade = 20;
float altura = 1.75;

if (idade > 18) {
    print(idade);
}
"""
\end{lstlisting}

\end{frame}

\begin{frame}{Resultado esperado}

    \begin{block}{Tokens identificados}
        \texttt{PALAVRA\_CHAVE: int}

        \texttt{IDENTIFICADOR: idade}

        \texttt{OPERADOR: =}

        \texttt{NUMINT: 20}

        \texttt{DELIMITADOR: ;}
    \end{block}

    \vspace{0.3cm}

    O processo é repetido até que todo o código-fonte
    seja analisado.

\end{frame}

% ------------------------------------------------
\section{Considerações finais}

\begin{frame}{Considerações finais}

    O desenvolvimento do analisador léxico permite
    compreender conceitos fundamentais da construção
    de linguagens de programação.

    \vspace{0.4cm}

    Além disso, estabelece uma base para a implementação
    de outras etapas de um compilador, como:

    \begin{itemize}
        \item Análise sintática;
        \item Construção de árvores sintáticas;
        \item Análise semântica;
        \item Geração de código intermediário.
    \end{itemize}

\end{frame}

% ------------------------------------------------
% ENCERRAMENTO
% ------------------------------------------------

\begin{frame}
    \centering

    {\Huge\color{primary} Obrigado!}

    \vspace{0.5cm}

    {\large Analisador Léxico -- MiniLang}

    \vspace{0.3cm}

    {\small Ciência da Computação}
\end{frame}

\end{document}